
This work sits at the intersection of LLM code evaluation, property-based testing for code verification, and formal contract checking.

\subsubsection{2.1 LLM Code Evaluation: From Sparse to Dense Testing}

HumanEval [Chen2021HumanEval] established functional correctness via unit tests as the primary evaluation axis for LLM-generated code. EvalPlus [Liu2023EvalPlus] addressed test-suite gaming by expanding HumanEval by $80\times$ and MBPP by $35\times$, reducing pass@1 by up to 23.1\%. This dense differential testing approach constitutes the current standard for test-based LLM code evaluation.

However, differential testing is limited to checking output equality on sampled inputs. It cannot verify whether a program satisfies properties that must hold for all valid inputs: relational invariants, quantified conditions, or semantic contracts specifying behavior beyond input-output pairs. This work builds on EvalPlus as the strongest available test-based baseline and quantifies the gap between it and a formal contract oracle.

\subsubsection{2.2 Property-Based Testing for LLM Code}

Bose [Bose2025Prompts] applied property-based testing (Hypothesis-style) to StarCoder and CodeLlama on MBPP and HumanEval, finding that 18--32\% of programs fail outright on property-based inputs. However, this work covers only two models, uses no formal benchmark-provided contract annotations, and lacks the oracle isolation design needed to separate input-generation advantage from oracle semantic strength.

Newcomb et al. [Newcomb2025Preconditions] examined how pre/post-condition constraints provided in prompts affect LLM code generation quality, showing that structural constraints improve correctness. That work studies contract-guided generation; this work studies contract-based evaluation of existing outputs.

No prior PBT work answers whether additional violations reflect oracle semantic strength or input exploration advantage. The oracle isolation design introduced in Section 3 resolves this directly.

\subsubsection{2.3 Formal Contract Checking for LLM Code}

ContractEval [Lim2025ContractEval] is the only publicly available benchmark pairing HumanEval+/MBPP+ tasks with formal Python pre/post-condition contracts. The original evaluation found 0\% contract satisfaction for 5 open-source models under standard prompting. Critically, ContractEval uses Z3 SMT checking, which is tractable for only 25.82\% of tasks. The execution-based approach introduced here achieves 100\% tractability across all 364 tasks and extends evaluation to closed-source models.

All claims in this related work survey are drawn from peer-reviewed or publicly archived sources.

\subsubsection{2.4 Summary}

\begin{table*}[htbp]
\centering
\resizebox{\dimexpr 2\columnwidth+\columnsep\relax}{!}{%
\begin{tabular}{lllll}
\toprule
Work & Oracle Type & Models & Oracle Isolation & Contract Annotations \\
\midrule
EvalPlus [Liu2023EvalPlus] & Differential (dense) & Multiple & --- & None \\
Bose [Bose2025Prompts] & PBT (no formal contracts) & 2 & No & Manual (researcher-authored) \\
ContractEval [Lim2025ContractEval] & SMT (25.82\% tractable) & 5 (open-source) & No & Formal \\
Newcomb et al. [Newcomb2025Preconditions] & Test-based & Multiple & No & Prompt-provided \\
**This work** & **Execution-based (100\% tractable)** & **5 (open+closed)** & **Yes (CVT design)** & **Formal** \\
\bottomrule
\end{tabular}
}
\end{table*}

